\documentclass[../../main.tex]{subfiles}% 注意这里的文件路径不能用 ./main.tex ,否则用latexmk编译子文件会报错
\graphicspath{{\subfix{./image/}}} % 指定图片目录,后续可以直接使用图片文件名
% 注意这里的文件路径不能用 ../../image/ ,否则用latexmk编译子文件会报错

% 例如:
% \begin{figure}[H]
% \centering
% \includegraphics[scale=0.3]{图.png}
% \caption{}
% \label{figure:图}
% \end{figure}
% 注意:上述\label{}一定要放在\caption{}之后,否则引用图片序号会只会显示??.

\begin{document}

\section{曲线的曲率和Frenet标架}

曲率刻画曲线的弯曲程度，是切线方向对弧长的变化率。
设$\bm{r}(s)$为以弧长$s$为参数的正则曲线，**单位切向量**为
\[
\bm{\alpha}(s) = \bm{r}'(s),\quad |\bm{\alpha}(s)|=1.
\]

\begin{definition}[曲率与曲率向量]
曲线的**曲率向量**为$\bm{\alpha}'(s)=\bm{r}''(s)$，**曲率**定义为
\[
\kappa(s) = \big|\bm{\alpha}'(s)\big| = \big|\bm{r}''(s)\big| \geq 0.
\]
\end{definition}

\begin{theorem}
曲线为**直线**的充要条件是其曲率$\kappa(s)\equiv 0$。
\end{theorem}


当$\kappa(s)\neq 0$时，定义**主法向量**：
\[
\bm{\beta}(s) = \dfrac{\bm{\alpha}'(s)}{\kappa(s)},\quad |\bm{\beta}(s)|=1,\ \bm{\beta}(s)\perp\bm{\alpha}(s).
\]

定义**次法向量**：
\[
\bm{\gamma}(s) = \bm{\alpha}(s) \times \bm{\beta}(s).
\]

\begin{definition}[Frenet标架]
由曲线的点$\bm{r}(s)$与单位正交右手标架$\{\bm{\alpha}(s),\bm{\beta}(s),\bm{\gamma}(s)\}$构成的标架
\[
\mathfrak{F}(s) = \big\{\bm{r}(s);\ \bm{\alpha}(s),\bm{\beta}(s),\bm{\gamma}(s)\big\}
\]
称为曲线在$s$处的**Frenet标架**，是曲线的内在标架。
\end{definition}


若曲线参数为$t$（非弧长），则曲率公式为
\[
\kappa(t) = \dfrac{\big|\bm{r}'(t)\times\bm{r}''(t)\big|}{\big|\bm{r}'(t)\big|^3}.
\]

\begin{example}
圆螺旋线$\bm{r}(t)=(a\cos t, a\sin t, bt)$（$a>0$）的曲率为常数：
\[
\kappa = \dfrac{a}{a^2+b^2}.
\]
\end{example}

\begin{exercise}
\begin{enumerate}
\item 求曲线$\bm{r}(t)=\big(at, a\sqrt{2}\ln t, a/t\big)$的曲率。
\item 求螺旋线$\bm{r}(t)=(a\cos t,a\sin t,bt)$的密切平面方程。
\end{enumerate}
\end{exercise}




\end{document}